%%%%%%%%%%%%%%%%%%%%%%%%%%%%%%%%%%%%%%%%%%%%%%%%%%%%%%%%%%%%%%%%%%%%%%
% Lớp 8 | Chương II. Hằng đẳng thức đáng nhớ và ứng dụng | Phân tích đa thức thành nhân tử
% Dạng 8D2?9-4: Ứng dụng: tìm x, tính nhanh, chứng minh chia hết
%%%%%%%%%%%%%%%%%%%%%%%%%%%%%%%%%%%%%%%%%%%%%%%%%%%%%%%%%%%%%%%%%%%%%%


% nguồn: De_AI_Toan8_cuoi_hoc_ki_i_1 (nhập 21/11/2026)
\begin{bt}%[8D2V9-4]%[ma:c3c26e16]
	Tìm $x$, biết:
a) $2x(x - 5) - 3(x - 5) = 0$;
b) $x^2 - 9 = 4x(x - 3)$;
c) $x^3 - 5x^2 + 6x = 0$.
	\loigiai{
		a) $2x(x - 5) - 3(x - 5) = 0$ \\
$(x - 5)(2x - 3) = 0$ \\
Suy ra $x - 5 = 0$ hoặc $2x - 3 = 0$, do đó $x = 5$ hoặc $x = \dfrac{3}{2}$. \\
b) $x^2 - 9 = 4x(x - 3)$ \\
$(x - 3)(x + 3) - 4x(x - 3) = 0$ \\
$(x - 3)(x + 3 - 4x) = 0$ \\
$(x - 3)(3 - 3x) = 0$ \\
Suy ra $x - 3 = 0$ hoặc $3 - 3x = 0$, do đó $x = 3$ hoặc $x = 1$. \\
c) $x^3 - 5x^2 + 6x = 0$ \\
$x(x^2 - 5x + 6) = 0$ \\
$x(x^2 - 2x - 3x + 6) = 0$ \\
$x(x - 2)(x - 3) = 0$ \\
Suy ra $x = 0$, $x = 2$ hoặc $x = 3$.
	}
\end{bt}

% nguồn: chuyen_dang (nhập 22/11/2026)
\begin{bt}%[8D2H9-4]%[ma:8fdb3b8e]
	Tìm $x$, biết $3x(x - 2) - x(3x + 4) = 20$.
	\loigiai{
		Ta có $3x(x - 2) - x(3x + 4) = 20$ \\
$3x^2 - 6x - 3x^2 - 4x = 20$ \\
$-10x = 20$ \\
$x = -2$. Vậy $x = -2$.
	}
\end{bt}

% nguồn: chuyen_dang (nhập 22/11/2026)
\begin{bt}%[8D2H9-4]%[ma:29e0aa94]
	Tìm $x$, biết $(2x - 1)(x + 3) - (2x + 5)(x - 2) = 15$.
	\loigiai{
		Ta có $(2x^2 + 6x - x - 3) - (2x^2 - 4x + 5x - 10) = 15$ \\
$(2x^2 + 5x - 3) - (2x^2 + x - 10) = 15$ \\
$4x + 7 = 15 \Rightarrow 4x = 8 \Rightarrow x = 2$. Vậy $x = 2$.
	}
\end{bt}

% nguồn: chuyen_dang (nhập 22/11/2026)
\begin{bt}%[8D2H9-4]%[ma:a27147c9]
	Tìm $x$, biết $(x + 2)(x - 5) - x(x + 1) = 10$.
	\loigiai{
		Ta có $(x^2 - 5x + 2x - 10) - (x^2 + x) = 10$ \\
$x^2 - 3x - 10 - x^2 - x = 10$ \\
$-4x - 10 = 10 \Rightarrow -4x = 20 \Rightarrow x = -5$. Vậy $x = -5$.
	}
\end{bt}

% nguồn: chuyen_dang (nhập 22/11/2026)
\begin{bt}%[8D2H9-4]%[ma:dab31780]
	Tìm $x$, biết $(3x - 2)(2x + 1) - 6x(x - 3) = 32$.
	\loigiai{
		Ta có $(6x^2 + 3x - 4x - 2) - (6x^2 - 18x) = 32$ \\
$6x^2 - x - 2 - 6x^2 + 18x = 32$ \\
$17x - 2 = 32 \Rightarrow 17x = 34 \Rightarrow x = 2$. Vậy $x = 2$.
	}
\end{bt}
